%% colophon.tex — production note (traditionally the last page).
%% Typeface names are dispatched on \BookFontProfile so the colophon
%% is always TRUE for the build that produced it — a colophon that
%% lies about its own fonts is the classic self-publishing tell.
\chapter{Colophon}

\noindent This book was designed and typeset with
\ifluatex Lua\LaTeX\else Xe\LaTeX\fi{} from a single-source template:
metadata, trim, typography, and edition structure all derive from one
configuration file, and every output format---print interior, bleed
variant, and ebook---builds from the same LaTeX sources.

\medskip

\noindent The text is set in
\ifdefstring{\BookFontProfile}{libertinus}{%
  Libertinus Serif, with Libertinus Sans for titling accents and
  Libertinus Mono for code%
}{}%
\ifdefstring{\BookFontProfile}{garamond}{%
  EB Garamond, with a companion humanist sans for titling accents and
  Noto Sans Mono for code%
}{}%
\ifdefstring{\BookFontProfile}{plex}{%
  IBM Plex Serif, with IBM Plex Sans for titling accents and IBM Plex
  Mono for code (SemiBold serving as bold throughout)%
}{}%
, at a base size of \BookBaseSize{} on a \BookTrimWidth{} by
\BookTrimHeight{} trim.

\medskip

\noindent Interior architecture: a modular preamble with a four-layer
color system, semantic callout boxes, and micro-typographic protrusion
tuned per typeface. Citations are managed with biblatex and biber.

\bigskip

\noindent\textit{\BookPublisher\ \textperiodcentered\ \BookYear}
